% Nutzeranalyse und Kontextanalyse

\subsubsection{Übersicht und Beschreibung von Benutzergruppen}

\smallskip
\noindent
Das Ziel der App ist es, Open Data der Stadt Wien für eine breite Öffentlichkeit zugänglich zu machen.
Basierend auf Nutzerinterviews und Recherchen lassen sich folgende Benutzergruppen identifizieren:

\medskip
\noindent\textbf{Primäre Nutzer:innen -- Gruppe 1: Technikaffine Studierende}

\smallskip
\noindent
Diese Gruppe ist technisch versiert und hat Erfahrung mit Daten und Visualisierungen.
Sie nutzen die App zur Recherche, für Uni-Projekte oder aus persönlichem Interesse.
Sie wünschen sich genaue, detaillierte Datensätze und aussagekräftige Visualisierungen, die auch mit anderen geteilt werden können.
\textit{Diese Gruppe wird durch die Persona Benjamin Eberhart repräsentiert.}
Potenzielle Probleme: Frustration bei zu vereinfachter Darstellung oder fehlenden Exportmöglichkeiten.

\medskip
\noindent\textbf{Primäre Nutzer:innen -- Gruppe 2: Content Creator und Blogger}

\smallskip
\noindent
Personen, die Daten als Inspirationsquelle oder Grundlage für eigene Inhalte verwenden.
Sie suchen nach visuell ansprechenden Darstellungen, die sie in Blogbeiträgen oder Social-Media-Posts einsetzen können.
Die App muss schnell verständlich sein und attraktive Visualisierungen liefern.
\textit{Diese Gruppe wird durch die Persona Leah Reiniger repräsentiert.}
Potenzielle Probleme: Zu komplexe Bedienung oder unattraktives Design könnte diese Gruppe abschrecken.

\medskip
\noindent\textbf{Primäre Nutzer:innen -- Gruppe 3: Berufliche Nutzer:innen}

\smallskip
\noindent
Diese Gruppe verwendet die App im beruflichen Kontext.
Sie benötigen spezifische Datensätze -- etwa zu Mietpreisen oder Wohnkosten -- und exportieren Visualisierungen für interne Berichte oder Präsentationen.
Sie schreiben keinen Code, erwarten aber eine hohe Datengenauigkeit und eine intuitive Oberfläche.
\textit{Diese Gruppe wird durch die Persona Ralph Ebersbach repräsentiert.}
Potenzielle Probleme: Unzureichende Exportfunktionen oder ungenaue Daten würden das Vertrauen dieser Gruppe stark beeinträchtigen.

\medskip
\noindent\textbf{Sekundäre Nutzer:innen -- Gelegenheitsnutzer:innen}

\smallskip
\noindent
Diese Gruppe hat kein vertieftes Interesse an Datenvisualisierung, nutzt die App aber dennoch,
um schnell zu interessanten Informationen zu gelangen -- etwa für Quizabende oder alltägliche Gespräche.
Die App muss intuitiv genug sein, damit diese Nutzer:innen ohne Vorkenntnisse schnell ans Ziel kommen.
\textit{Diese Gruppe wird durch die Persona Luca Kunze repräsentiert.}

\medskip
\noindent\textbf{Negative Persona -- Problematische Nutzer:innen}

\smallskip
\noindent
Daten können nicht nur für positive Zwecke verwendet werden.
Diese Persona repräsentiert Nutzer:innen, die versuchen, Datenzusammenhänge zu missbrauchen -- etwa um manipulative oder irreführende Inhalte zu erstellen.
Sie sind technisch versiert und wünschen sich möglichst offenen, anonymen Zugang zu so vielen Daten wie möglich.
Für die App ergeben sich daraus wichtige Designanforderungen: Persönliche Daten müssen anonymisiert und geschützt werden.
Alle bereitgestellten Daten müssen auf Richtigkeit überprüft sein und aus vertrauenswürdigen, zuverlässigen Quellen stammen.

\subsubsection{Aufgaben und Ziele der Nutzer:innen}

\smallskip
\noindent
Die primären Aufgaben der Nutzer:innen lassen sich gruppenübergreifend in folgende Kategorien einteilen:

\begin{itemize}
    \item \textbf{Erkunden und Entdecken:} Nutzer:innen wie Benjamin möchten neue, unbekannte Datensätze der Stadt Wien finden und durchstöbern -- ohne ein konkretes Ziel vorab zu haben.
    \item \textbf{Gezielt Suchen:} Nutzer:innen wie Ralph haben ein spezifisches Informationsbedürfnis (z.\,B. Mietpreise, Bevölkerungsdaten) und möchten schnell zum richtigen Datensatz gelangen.
    \item \textbf{Visualisieren und Teilen:} Nutzer:innen wie Leah möchten Daten in ansprechenden Visualisierungen darstellen und diese einfach exportieren oder in sozialen Netzwerken teilen.
    \item \textbf{Schnell informieren:} Nutzer:innen wie Luca wollen ohne Vorkenntnisse unkompliziert zu einer interessanten Information gelangen.
\end{itemize}

\subsubsection{Potenzielle Probleme mit dem System}

\begin{itemize}
    \item \textbf{Zu komplexe Navigation:} Nutzer:innen ohne technischen Hintergrund könnten überfordert sein.
    \item \textbf{Lange Ladezeiten:} Bei mobiler Nutzung mit schwacher Verbindung kritisch.
    \item \textbf{Unklare Datenquellen:} Nutzer:innen könnten der Richtigkeit der Daten misstrauen.
    \item \textbf{Fehlende Exportfunktion:} Besonders für berufliche Nutzer:innen ein Problem.
    \item \textbf{Datenmissbrauch:} Ohne Anonymisierung könnten sensible Daten missbraucht werden.
\end{itemize}

\subsubsection{Kontextanalyse}

\smallskip
\noindent
Die Anwendung wird primär als mobile App auf Smartphones genutzt. Daraus ergeben sich folgende Nutzungskontexte:

\begin{itemize}
    \item \textbf{Nutzungsort und -situation:} Die App wird häufig unterwegs verwendet -- in öffentlichen Verkehrsmitteln, in Pausen oder zuhause. Es ist daher mit kurzen, spontanen Nutzungssessions zu rechnen. Berufliche Nutzer:innen wie Ralph verwenden die App hingegen auch am Arbeitsplatz für gezielte Recherchen.
    \item \textbf{Technische Rahmenbedingungen:} Da die Nutzung primär mobil erfolgt, muss die App auch bei schwacher Internetverbindung grundlegend funktionieren. Ladezeiten sollten minimal sein. Die begrenzte Bildschirmgröße erfordert kompakte, dennoch gut lesbare Visualisierungen.
    \item \textbf{Zeitlicher Kontext:} Die meisten Nutzer:innen öffnen die App mit einem konkreten Informationsbedürfnis. Die Navigation zu den gewünschten Daten muss daher schnell und effizient sein -- komplexe Einstiegshürden sind zu vermeiden.
    \item \textbf{Soziale Aspekte:} Ein Teil der Nutzer:innen teilt Inhalte aus der App in sozialen Netzwerken oder im persönlichen Gespräch. Eine einfache Teilen- und Exportfunktion würde den Mehrwert der App für alle Gruppen deutlich erhöhen.
\end{itemize}
